
\noindent {\bfseries GFT - Fundamentos de Cloud com AWS} \hfill\ Jun 2026 -- Jul 2026 \\
Bootcamp da DIO, para aprender sobre a AWS, de forma prática.
\begin{itemize}
    \item Criei contas e configurei cargos e permissões com o IAM;
    \item Criei diagramas  e fluxogramas de arquitetura e infraestrutura nuvem AWS no draw.io;
    \item Estudei e apliquei sobre Amazon EC2, Lambda, Serverless;
    \item Estudei e apliquei sobre os bancos de dados RDS, DynamoDB, S3;
    \item Estudei e apliquei sobre SQS e SNS;
    \item Aprimorei meus conhecimentos sobre EKS, Docker e DevOPS.
\end{itemize}
\noindent {\bfseries Tecnologias:} AWS, Cloud, Amazon, EC2, Lambda, IAM, RDS, DynamoDB, S3, SQS, SNS, EKS, DevOps.

%-----------------------------------------    
        
\noindent {\bfseries Estudos sobre TDD} \hfill\ Mai 2026 \\
Estudos sobre Test-Driven Development, técnica de desenvolvimento que utiliza testes de unidade como guias do código.
\begin{itemize}
    \item Aprendi sobre TDD, o que é, que ele mira nos testes de unidade;
    \item Entendi que ele foca no resultado esperado, e não como chegar no resultado;
    \item Entendi a ideia do Test First, Red/Refactor/Green e Baby Steps, passo por passo, onde testamos primeiro;
    \item Aprendi que o TDD fala que não devemos fazer que os códigos funcionem de primeira, precisamos refatorar.
\end{itemize}
\noindent {\bfseries Tecnologias:} TDD, Domain, Value Objects, Context.

%-----------------------------------------    
        
\noindent {\bfseries Desafio DDD: Sistema CineFlow} \hfill\ Mai 2026 \\
Exercício prático de modelagem de um sistema de gerenciamento de salas de cinema aplicando os conceitos de Domain-Driven Design.
\begin{itemize}
    \item Interpretei o cenário e identifiquei as ações e entidades do sistema;
    \item Refleti e defini os domínios e seus tipos como principal, genérico e auxiliar;
    \item Mapeei os limites de cada domínio e o que pertence a cada contexto delimitado;
    \item Defini os contextos e refleti sobre microsserviços com comunicação assíncrona via RabbitMQ;
    \item Identifiquei e estruturei Value Objects como assentos e valores monetários.
\end{itemize}
\noindent {\bfseries Tecnologias:} DDD, Domain, Value Objects, Context.

%-----------------------------------------    
        
\noindent {\bfseries Estudos sobre DDD} \hfill\ Abr 2026 -- Mai 2026 \\
Estudos sobre Domain-Driven Design focados na modelagem de domínios complexos, limites de contextos e arquitetura de software.
\begin{itemize}
    \item Entendi que o Domínio é o coração do negócio;
    \item Compreendi os tipos de domínios divididos em principal, genérico e auxiliar;
    \item Assimilei a importância da Linguagem Ubíqua e da comunicação com o Domain Expert;
    \item Aprendi a delimitar os contextos da aplicação e a estabelecer as fronteiras de cada um.
\end{itemize}
\noindent {\bfseries Tecnologias:} DDD, Domain, Value Objects, Context.

%-----------------------------------------    
        
\noindent {\bfseries RabbitMQ: Queues, Exchanges e DLQ} \textbar{}  \href{https://github.com/Lucas-Pontes-Soares/learn-rabbitmq}{github.com/Lucas-Pontes-Soares/learn-rabbitmq} \hfill\ Abr 2026 \\
Estudo prático de mensageria em uma arquitetura de microsserviços de e-commerce com foco em resiliência e tratamento de erros.
\begin{itemize}
    \item Estudei sobre os conceitos fundamentais do RabbitMQ;
    \item Desenvolvi os microsserviços de forma independente e especializada;
    \item Desenvolvi o publicador e os consumidores de mensagens do ecossistema;
    \item Implementei a lógica de Dead Letter Queue para tratamento de mensagens com falha;
    \item Utilizei Exchanges estruturadas para garantir o roteamento correto e indireto para as filas.
\end{itemize}
\noindent {\bfseries Tecnologias:} RabbitMQ, Node.js, Docker, Queues, Exchanges, DLQ.

%-----------------------------------------    
        
\noindent {\bfseries Submissão Artigo na Revista InterAgro} \textbar{}  \href{https://github.com/Lucas-Pontes-Soares/estacao-agroclimatica-inteligente}{github.com/Lucas-Pontes-Soares/estacao-agroclimatica-inteligente} \hfill\ Set 2025 -- Abr 2026 \\
Estação agroclimática inteligente baseada em IoT e protocolo MQTT utilizando a placa ESP32 para monitoramento ambiental.
\begin{itemize}
    \item Estudei sobre o funcionamento e aplicação do protocolo MQTT;
    \item Desenvolvi e modelei a estrutura física simulada do ESP32 na plataforma Wokwi;
    \item Programei a lógica de captura de dados de sensores e mensageria;
    \item Configurei o envio e o recebimento de notificações e alertas estruturados via broker MQTT.
\end{itemize}
\noindent {\bfseries Tecnologias:} ESP32, Python, Arduino, Agro, Revista, InterAgro, Wokwi.

%-----------------------------------------    
        
\noindent {\bfseries Back-end Node com Testes Unitários, Integração e E2E} \textbar{}  \href{https://github.com/Lucas-Pontes-Soares/node-testable-apps}{github.com/Lucas-Pontes-Soares/node-testable-apps} \hfill\ Abr 2026 \\
API de agendamentos para uma barbearia desenvolvida de forma isolada do banco de dados utilizando os princípios do SOLID.
\begin{itemize}
    \item Estudei sobre conceitos de Testes Unitários, de Integração e E2E;
    \item Utilizei conceitos do SOLID, inversão de dependências e POO;
    \item Construi um serviço independente de banco de dados, criando uma aplicação isolada;
    \item Desenvolvi testes unitários com Vitest no Node.js para validar regras de negócio e datas.
\end{itemize}
\noindent {\bfseries Tecnologias:} Node.js, Typescript, Git, Vitest, Tests, Unit Tests, Integration Tests, E2E Tests, SOLID, POO.

%-----------------------------------------    
        
\noindent {\bfseries NestJS Library API} \textbar{}  \href{https://github.com/Lucas-Pontes-Soares/nest-library-api}{github.com/Lucas-Pontes-Soares/nest-library-api} \hfill\ Mar 2026 -- Abr 2026 \\
API REST para gerenciamento de uma biblioteca física, implementando controle de acesso por perfis (RBAC) e regras de permissões.
\begin{itemize}
    \item Apliquei práticas de arquitetura modular separando recursos em repositórios, DTOs, controllers, interfaces e services;
    \item Implementei autenticação via tokens JWT utilizando Guards do NestJS para a validação das requisições;
    \item Desenvolvi a camada de persistência com banco de dados PostgreSQL utilizando o Drizzle ORM;
    \item Desenvolvi a lógica de controle de acesso baseado em funções (RBAC) separando permissões entre clientes e bibliotecários;
    \item Estruturei uma matriz de permissões no arquivo permissions.ts para gerenciar dinamicamente as ações permitidas a cada perfil.
\end{itemize}
\noindent {\bfseries Tecnologias:} NestJS, Node.js, Typescript, Git, PostgreSQL, Drizzle, Roles, RBAC, ABAC, JWT.

%-----------------------------------------    
        
\noindent {\bfseries DevRoast} \textbar{}  \href{https://github.com/Lucas-Pontes-Soares/devroast}{github.com/Lucas-Pontes-Soares/devroast} \hfill\ Mar 2026 \\
Analisador de qualidade de código baseado em inteligência artificial que avalia e fornece feedbacks de 0 a 10.
\begin{itemize}
    \item Apliquei melhores práticas de engenharia de prompts estruturados para guiar os agentes de IA;
    \item Integrei o Claude Code com Model Context Protocol (MCP) e skills customizadas para automação do fluxo;
    \item Apliquei a metodologia Spec-Driven Development (SSD) e etapas de plan-build;
    \item Desenvolvi a interface web responsiva utilizando a arquitetura do Next.js;
    \item Desenvolvi a camada de persistência e serviços back-end e banco de dados PostgreSQL;
    \item Implementei a lógica de análise estática e interpretativa de código conectando à API do Gemini.
\end{itemize}
\noindent {\bfseries Tecnologias:} Next.js, IA, Gemini, Claude Code, MCP, Git, Agentes, Typescript, tailwind-css, Pencil, Base-ui, PostgreSQL.

%-----------------------------------------    
        
\noindent {\bfseries MasterClass NestJS} \textbar{}  \href{https://github.com/Lucas-Pontes-Soares/masterclass-nestjs}{github.com/Lucas-Pontes-Soares/masterclass-nestjs} \hfill\ Mar 2026 \\
Estudo das práticas e arquitetura opinada do framework NestJS, explorando injeção de dependência e desacoplamento de camadas.
\begin{itemize}
    \item Desenvolvi uma API REST em NestJS para absorver a arquitetura e o ecossistema do framework;
    \item Apliquei na prática os conceitos de Inversão de Dependência e Injeção de Dependência para desacoplamento;
    \item Configurei e integrei um banco de dados SQLite utilizando o Prisma ORM para persistência dos dados;
    \item Organizei e estruturei a aplicação dividindo as responsabilidades em módulos, controllers, providers, repositories e DTOs.
\end{itemize}
\noindent {\bfseries Tecnologias:} NestJS, Typescript, Node.js, Prisma, SQLite, Git.

%-----------------------------------------    
        
\noindent {\bfseries Práticas de RAG} \textbar{}  \href{https://github.com/Lucas-Pontes-Soares/learn-rag}{github.com/Lucas-Pontes-Soares/learn-rag} \hfill\ Mar 2026 \\
Estudo prático de Retrieval-Augmented Generation abordando chunking, embeddings, banco de dados vetorial e busca por similaridade.
\begin{itemize}
    \item Desenvolvi lógica para extrair conteúdo em formato Markdown de sites utilizando a API do Firecrawl;
    \item Implementei a segmentação de documentos em chunks aplicando estratégias de overlap para preservação de contexto;
    \item Gerei vetores de embeddings para cada fragmento de texto utilizando os modelos da API do Gemini;
    \item Persisti os embeddings gerados estruturadamente no banco de dados vetorial Turso;
    \item Converteri as consultas textuais dos usuários em embeddings para viabilizar a indexação;
    \item Executei algoritmos de busca por similaridade vetorial para extrair os contextos mais relevantes para a IA.
\end{itemize}
\noindent {\bfseries Tecnologias:} Typescript, RAG, IA, Docker, Turso, Firecrawl, Gemini, Git, Embeddings, Chunks.

%-----------------------------------------    
        
\noindent {\bfseries Cache com Redis em Node.js} \textbar{}  \href{https://github.com/Lucas-Pontes-Soares/learn-redis}{github.com/Lucas-Pontes-Soares/learn-redis} \hfill\ Mar 2026 \\
Estudo e implementação de estratégias de cache do tipo chave-valor utilizando Redis em uma API Node.js para otimização.
\begin{itemize}
    \item Desenvolvi o fluxo completo para criação e persistência de dados de usuários;
    \item Configurei o ambiente em contêineres Docker para gerenciar o banco de dados e a instância do Redis;
    \item Configurei o Redis como camada de cacheamento para otimizar o acesso aos dados cadastrados;
    \item Implementei o padrão de chaves estruturadas serializando objetos complexos de usuários em strings;
    \item Analisei as diferenças arquiteturais e estruturais entre bancos de dados relacionais e o Redis;
    \item Avaliei o impacto na performance mensurando a redução no tempo de resposta das requisições.
\end{itemize}
\noindent {\bfseries Tecnologias:} Node.js, Redis, PostgreSQL, Git, Docker.

%-----------------------------------------    
        
\noindent {\bfseries Desafio Microsserviços Escaláveis Node.js} \textbar{}  \href{https://github.com/Lucas-Pontes-Soares/desafio-microsservicos-escalaveis-nodejs}{github.com/Lucas-Pontes-Soares/desafio-microsservicos-escalaveis-nodejs} \hfill\ Fev 2026 \\
Desenvolvimento de ecossistema baseado em microsserviços independentes, focando em alta disponibilidade de maneira escalável.
\begin{itemize}
    \item Desenvolvi microsserviços independentes com comunicação assíncrona via RabbitMQ;
    \item Configurei o ambiente Docker para a execução local dos bancos de dados e do message broker;
    \item Apliquei a camada de persistência utilizando PostgreSQL integrado ao Drizzle ORM;
    \item Implementei e gerenciei o fluxo de migrations automatizadas no banco de dados;
    \item Implementei a estrutura de observabilidade da arquitetura utilizando Grafana e Jaeger para logging e tracing;
    \item Automatizei o fluxo de deploy em infraestrutura AWS Fargate utilizando a ferramenta de IaC Pulumi.
\end{itemize}
\noindent {\bfseries Tecnologias:} Node.js, PostgreSQL, RabbitMQ, Grafana, Jaeger, Pulumi, AWS, Git, GitHub, Drizzle, Docker, GitHub Actions, CI, CD.

%-----------------------------------------    
        
\noindent {\bfseries Dumble: Sistema Educacional Com Inteligência Artificial} \textbar{}  \href{https://github.com/Lucas-Pontes-Soares/dumble}{github.com/Lucas-Pontes-Soares/dumble} \hfill\ Fev 2025 -- Dez 2025 \\
Projeto integrador da faculdade, uma plataforma web interativa que centraliza o ensino para fortalecer o vínculo professor-aluno.
\begin{itemize}
    \item Desenvolvi a parte da IA com RAG, onde os professores anexavam arquivos e o conteudo deles eram disponibilizados para a IA;
    \item Construi os prompts rigorosos e estruturados de IA, para o chat-bot e a sugestão de perguntas;
    \item Desenvolvi o front-end de maneira responsiva;
    \item Implementei autenticação de usuários via tokens JWT;
    \item Gerenciei o banco de dados SQL;
    \item Administrei e realizei o deploy (hospedagem).
\end{itemize}
\noindent {\bfseries Tecnologias:} IA, Node.js, React.js, PostgreSQL (SQL), RAG, JWT, Heroku, Git, tailwind-css, shadcn-ui.

%-----------------------------------------    
        
\noindent {\bfseries Janela Automática} \textbar{}  \href{https://github.com/Lucas-Pontes-Soares/JanelaAutomatica}{github.com/Lucas-Pontes-Soares/JanelaAutomatica} \hfill\ Set 2025 -- Dez 2025 \\
Sistema de automação residencial utilizando ESP32, sensores ambientais e protocolo MQTT para o controle inteligente de janela.
\begin{itemize}
    \item Desenvolvi o firmware para o microcontrolador ESP32, gerenciando a leitura de sensores de chuva, temperatura e luminosidade;
    \item Projetei a lógica de automação segmentada para controlar de forma independente o vidro, a persiana e a veneziana;
    \item Implementei a telemetria e a comunicação dos dados via protocolo MQTT para integração com um aplicativo externo.
\end{itemize}
\noindent {\bfseries Tecnologias:} Arduino, ESP-32, MQTT, Python.

%-----------------------------------------    
        
\noindent {\bfseries Avanti Intelligence: IA para detectar pneumonia} \textbar{}  \href{https://github.com/Lucas-Pontes-Soares/avanti-intelligence}{github.com/Lucas-Pontes-Soares/avanti-intelligence} \hfill\ Mai 2025 -- Out 2025 \\
Projeto para apresentar no evento da Geniuscon 2025 em Jacarezinho, conseguimos chegar nas finais.
\begin{itemize}
    \item Desenvolvemos e estruturamos o modelo de IA realizando o treinamento com um dataset real de imagens de raio-X de tórax;
    \item Realizamos o treinamento supervisionado da Rede Neural Convolucional (CNN);
    \item Analisamos as métricas obtidas e alcançamos uma taxa de acerto de 90% na identificação positiva da doença;
    \item Implementei o fluxo de autenticação de usuários na plataforma web utilizando tokens JWT;
    \item Aplicamos o pipeline de dados segmentando o conjunto de imagens em etapas de treinamento, validação e teste;
    \item Desenvolvi a interface web para apresentação do projeto e integrei o modelo exportado para a execução de testes em tempo real.
\end{itemize}
\noindent {\bfseries Tecnologias:} IA, Python, CNN, React.js, Aprendizado Supervisionado, Git, tailwind-css, shadcn-ui.

%-----------------------------------------    
        
\noindent {\bfseries Projeto Kubernetes Azure} \textbar{}  \href{https://github.com/Lucas-Pontes-Soares/atividade-interdisciplinar-kubernetes}{github.com/Lucas-Pontes-Soares/atividade-interdisciplinar-kubernetes} \hfill\ Out 2025 \\
Projeto de DevOps e Cloud focada em conteinerização com Docker e orquestração com Kubernetes na Azure.
\begin{itemize}
    \item Versionei o código-fonte da aplicação utilizando o ecossistema Git e GitHub;
    \item Configurei o Dockerfile para empacotar o site estático integrado ao servidor Nginx;
    \item Implementei pipelines de CI com GitHub Actions para validação e linter automatizado de HTML e CSS;
    \item Configurei o cluster Kubernetes na Azure (AKS) definindo a execução resiliente de dois pods;
    \item Realizei o deploy e a sustentação da aplicação conteinerizada utilizando a infraestrutura de nuvem da Azure.
\end{itemize}
\noindent {\bfseries Tecnologias:} Docker, DockerHub, CI, CD, GitHub, Git, GitHub Actions, Azure, Kubernetes.

%-----------------------------------------    
        
\noindent {\bfseries Maratona de Programação} \hfill\ Mar 2024 -- Out 2025 \\
Competições acadêmicas de programação (InterFatecs), resolução de problemas complexos de lógica e algoritmos.
\begin{itemize}
    \item Atuei em equipe compartilhando estratégias e dividindo a resolução de problemas sob pressão de tempo;
    \item Desenvolvi algoritmos otimizados e estruturas de dados eficientes na linguagem C para atender aos limites de execução;
    \item Realizei a análise dos enunciados complexos para extração de requisitos, cenários de teste e modelagem inicial das soluções.
\end{itemize}
\noindent {\bfseries Tecnologias:} C, Lógica, Maratona.

%-----------------------------------------    
        
\noindent {\bfseries Jogo da Velha com IA} \textbar{}  \href{https://github.com/Lucas-Pontes-Soares/jogo-da-velha-com-ia}{github.com/Lucas-Pontes-Soares/jogo-da-velha-com-ia} \hfill\ Set 2025 \\
Um Jogo da Velha clássico desenvolvido em Python onde o adversário é uma inteligência artificial integrada à API do Gemini.
\begin{itemize}
    \item Desenvolvi a lógica do jogo da velha gerenciando as rodadas entre o jogador e a IA;
    \item Implementei a funcionalidade de escolha de símbolo para o jogador entre as opções 'X' ou 'O';
    \item Apliquei níveis de dificuldade dinâmicos alterando as instruções de sistema enviadas à IA;
    \item Programei um sorteio aleatório com probabilidade equilibrada para definir o responsável pela jogada inicial;
    \item Estruturei o tabuleiro 3x3 através de uma matriz indexada para controle preciso das linhas e colunas;
    \item Implementei algoritmos de validação para verificação de condições de vitória, derrota ou empate.
\end{itemize}
\noindent {\bfseries Tecnologias:} Python, IA, Gemini, Git.

%-----------------------------------------    
        
\noindent {\bfseries Jornal Docker} \textbar{}  \href{https://github.com/Lucas-Pontes-Soares/mini-site-docker}{github.com/Lucas-Pontes-Soares/mini-site-docker} \hfill\ Set 2025 \\
Site temático de jornal para explicar sobre Docker utilizado DevOps e CI/CD.
\begin{itemize}
    \item Desenvolvi a interface front-end contendo páginas informativas sobre o ecossistema Docker;
    \item Implementei pipelines de Integração Contínua (CI) utilizando GitHub Actions;
    \item Configurei o Dockerfile para build de produção e o Docker Compose para otimizar o ambiente de desenvolvimento;
    \item Automatizei gatilhos para que a cada push ou pull request a aplicação fosse construída e testada de forma consistente;
    \item Realizei a hospedagem e a publicação automatizada do site através do GitHub Pages.
\end{itemize}
\noindent {\bfseries Tecnologias:} Docker, build, HTML, CSS, CI, GitHub, Git, GitHub Actions, GitHub Pages.

%-----------------------------------------    
        
\noindent {\bfseries WineQuality: Machine Learning, Árvore de Decisão para Classificar Vinhos} \textbar{}  \href{https://github.com/Lucas-Pontes-Soares/wine-quality}{github.com/Lucas-Pontes-Soares/wine-quality} \hfill\ Jun 2025 \\
Algoritmo de Árvore de Decisão desenvolvido para classificar a qualidade de vinhos com base em atributos químicos e físicos.
\begin{itemize}
    \item Treinamos o modelo preditivo utilizando um conjunto de dados supervisionado e rotulado;
    \item Implementamos e estruturamos algoritmos de Árvore de Decisão para a classificação dos dados;
    \item Executamos o pipeline completo de Machine Learning englobando as etapas de treinamento, validação, teste e avaliação métrica.
\end{itemize}
\noindent {\bfseries Tecnologias:} Python, IA, Árvore de Decisão.

%-----------------------------------------    
        
\noindent {\bfseries MazeSolver: Resolução de Labirinto com Busca Cega} \textbar{}  \href{https://github.com/Lucas-Pontes-Soares/mazeSolver}{github.com/Lucas-Pontes-Soares/mazeSolver} \hfill\ Mai 2025 \\
Algoritmo em Python que gera labirintos aleatórios e aplica a busca em largura (BFS) para encontrar o caminho entre dois pontos.
\begin{itemize}
    \item Desenvolvemos um algoritmo baseado no método de busca cega em largura (BFS) para exploração de grafos;
    \item Desenvolvemos uma lógica estruturada para gerar labirintos aleatórios com garantia de caminhos possíveis;
    \item Implementamos uma interface de visualização em tempo real para exibir o labirinto sendo descoberto pelas iterações;
    \item Analisamos e observamos na prática o comportamento, a ordem de expansão de nós e a eficiência da busca em largura.
\end{itemize}
\noindent {\bfseries Tecnologias:} Python, Busca cega, Navegação.

%-----------------------------------------    
        
\noindent {\bfseries Games Spring} \textbar{}  \href{https://github.com/Lucas-Pontes-Soares/games-spring}{github.com/Lucas-Pontes-Soares/games-spring} \hfill\ Jan 2025 \\
Projeto em Java Spring Boot integrado com front-end em JSP para o gerenciamento e CRUD de jogos, categorias e plataformas.
\begin{itemize}
    \item Utilizei docker-compose para subir o banco de dados SQL;
    \item Trabalhei no Spring Boot com Controllers, Models e Repositories;
    \item Construi a interface com JSP com o CRUD completo do jogo, categoria e plataforma;
    \item Trabalhei com rotas e metodos HTTP, parámetros e retornos na API.
\end{itemize}
\noindent {\bfseries Tecnologias:} Java, Spring Boot, Gradle, JSP, SQL, Docker, HTML, CSS.

%-----------------------------------------    
        
\noindent {\bfseries Sistema para Assistente Social} \hfill\ Fev 2024 -- Dez 2024 \\
Projeto acadêmico de engenharia de software focado em análise, modelagem do banco de dados para a assistência social.
\begin{itemize}
    \item Idealizamos o escopo funcional e as regras de negócio de todo o ecossistema do projeto;
    \item Modelamos a arquitetura do sistema através de diagramas de casos de uso, classes, atividades e dicionário de dados;
    \item Construímos protótipos de alta fidelidade para validação de interfaces e experiência do usuário;
    \item Desenvolvemos a estrutura do banco de dados relacional em SQL com tabelas, restrições e relacionamentos;
    \item Realizamos o processo de ETL para migração de dados históricos legados a partir de planilhas eletrônicas.
\end{itemize}
\noindent {\bfseries Tecnologias:} Banco de Dados, SQL, Análise de Requisitos.

%-----------------------------------------    
        
\noindent {\bfseries CETAF} \textbar{}  \href{https://github.com/Lucas-Pontes-Soares/CETAF}{github.com/Lucas-Pontes-Soares/CETAF} \hfill\ Out 2024 -- Dez 2024 \\
Sistema de gerenciamento de matrículas acadêmicas desenvolvido sob metodologias ágeis.
\begin{itemize}
    \item Desenvolvi a infraestrutura e os serviços do back-end utilizando o ecossistema Node.js;
    \item Gerenciei e estruturei o banco de dados relacional utilizando o SGDB MySQL;
    \item Construí e otimizei queries SQL nativas para manipulação e junção de dados complexos na camada de persistência.
\end{itemize}
\noindent {\bfseries Tecnologias:} React.js, Node.js, API, SQL, MySQL.

%-----------------------------------------    
        
\noindent {\bfseries Lista Personagens} \textbar{}  \href{https://github.com/Lucas-Pontes-Soares/ListagemPersonagens}{github.com/Lucas-Pontes-Soares/ListagemPersonagens} \hfill\ Out 2024 -- Dez 2024 \\
Aplicativo mobile desenvolvido em Flutter para o gerenciamento, listagem e detalhamento de personagens integrados a uma API.
\begin{itemize}
    \item Configurei e estruturei uma API REST simulada utilizando a plataforma MockAPI;
    \item Modelei o esquema de dados e os atributos fundamentais para a representação dos personagens;
    \item Implementei endpoints para a listagem completa, paginação e filtragem dinâmica dos registros;
    \item Desenvolvi a interface do aplicativo mobile em Flutter utilizando o FlutLab para consumo e renderização dos dados da API.
\end{itemize}
\noindent {\bfseries Tecnologias:} Flutter, Front-end, API.

%-----------------------------------------    
        
\noindent {\bfseries Ministrei mini curso de Robocode} \hfill\ Out 2024 \\
Organização e docência de um minicurso prático sobre a plataforma Robocode, culminando em um campeonato de robôs virtuais.
\begin{itemize}
    \item Apresentei os conceitos fundamentais e a arquitetura de simulação da plataforma Robocode para os alunos;
    \item Desenvolvi um robô exemplo em tempo real, explicando a lógica de programação, movimentação e escaneamento de alvos;
    \item Orientei os estudantes na prática em relação à API da ferramenta e estratégias de tomada de decisão baseada em eventos;
    \item Organizei e gerenciei o chaveamento de um campeonato competitivo de combate entre os tanques customizados criados pelos alunos.
\end{itemize}
\noindent {\bfseries Tecnologias:} Arduino, Robocode, Minicurso.

%-----------------------------------------    
        
\noindent {\bfseries Notes} \textbar{}  \href{https://github.com/Lucas-Pontes-Soares/NLW-Expert}{github.com/Lucas-Pontes-Soares/NLW-Expert} \hfill\ Fev 2024 \\
Aplicação web de gerenciamento de notas com suporte a criação de registros via texto ou por reconhecimento de voz.
\begin{itemize}
    \item Desenvolvi a interface front-end baseada em componentes reativos utilizando ReactJS e Typescript no ecossistema Vite;
    \item Implementei a estilização utilitária e totalmente responsiva do layout através das classes do Tailwind CSS;
    \item Integrei a SpeechRecognition API para viabilizar a criação automatizada de notas a partir da transcrição de voz do usuário;
    \item Estruturei as funcionalidades de indexação, busca dinâmica por palavras-chave e exclusão de registros na aplicação.
\end{itemize}
\noindent {\bfseries Tecnologias:} React.js, Front-end, Typescript.

%-----------------------------------------    
        
\noindent {\bfseries NewsBlog} \textbar{}  \href{https://github.com/Lucas-Pontes-Soares/NewsBlog}{github.com/Lucas-Pontes-Soares/NewsBlog} \hfill\ Fev 2024 \\
Plataforma de blog de artigos desenvolvida em ReactJS com suporte à criação e leitura de publicações estruturadas.
\begin{itemize}
    \item Desenvolvi a interface front-end baseada em componentes utilizando ReactJS, Typescript e a arquitetura do Vite;
    \item Implementei a persistência local do catálogo de postagens utilizando a API do LocalStorage;
    \item Estruturei a área de criação de artigos contendo formulários dinâmicos para título, subtítulo, imagem de capa e corpo de conteúdo;
    \item Desenvolvi efeitos de interação (hover) e gerenciamento de rotas para exibição e leitura completa das publicações.
\end{itemize}
\noindent {\bfseries Tecnologias:} React.js, Front-end, Typescript.

%-----------------------------------------    
        
\noindent {\bfseries Lista de Compras} \textbar{}  \href{https://github.com/Lucas-Pontes-Soares/ListaCompras}{github.com/Lucas-Pontes-Soares/ListaCompras} \hfill\ Jan 2024 \\
Aplicação web de lista de compras desenvolvida em Angular com operações de CRUD e persistência no LocalStorage.
\begin{itemize}
    \item Desenvolvi a interface front-end utilizando componentes e a estrutura arquitetural do ecossistema Angular;
    \item Implementei o ciclo de vida dos componentes para gerenciar o estado e a reatividade da aplicação;
    \item Configurei a persistência da lista de compras localmente utilizando a API do LocalStorage;
    \item Estruturei a lógica de controle de fluxos para gerenciamento, edição e exclusão de itens cadastrados.
\end{itemize}
\noindent {\bfseries Tecnologias:} Angular, Front-end.

%-----------------------------------------    
        
\noindent {\bfseries Memoteca} \textbar{}  \href{https://github.com/Lucas-Pontes-Soares/memoteca}{github.com/Lucas-Pontes-Soares/memoteca} \hfill\ Jan 2024 \\
Aplicação web desenvolvida em Angular 14 para o gerenciamento e CRUD completo de pensamentos, citações e trechos de músicas.
\begin{itemize}
    \item Desenvolvi a interface front-end modularizada utilizando a arquitetura de componentes do Angular 14;
    \item Implementei serviços para a comunicação estruturada e manipulação de estado dos dados da aplicação;
    \item Configurei a persistência e o armazenamento local de registros utilizando a API de LocalStorage;
    \item Estruturei as rotas e a lógica de validação para o fluxo de criação, leitura, edição e exclusão de cartões.
\end{itemize}
\noindent {\bfseries Tecnologias:} Angular, Front-end.

%-----------------------------------------    
        
\noindent {\bfseries AluraBooks} \textbar{}  \href{https://github.com/Lucas-Pontes-Soares/alurabooks}{github.com/Lucas-Pontes-Soares/alurabooks} \hfill\ Jan 2024 \\
Aplicação web desenvolvida em Angular focada em formulários orientados a templates com validações e integração com API de CEP.
\begin{itemize}
    \item Desenvolvi a interface front-end utilizando componentes estruturados e estilização responsiva no ecossistema Angular;
    \item Implementei formulários orientados a templates aplicando diretivas customizadas e restrições de validação de campos;
    \item Integrei uma API externa de busca de endereços para realizar o autopreenchimento dos campos a partir do CEP informado.
\end{itemize}
\noindent {\bfseries Tecnologias:} Angular, Front-end.

%-----------------------------------------    
        
\noindent {\bfseries Java com Spring Boot} \textbar{}  \href{https://github.com/Lucas-Pontes-Soares/todolist}{github.com/Lucas-Pontes-Soares/todolist} \hfill\ Out 2023 \\
API REST de uma aplicação To-Do List desenvolvida para consolidar conceitos iniciais de rotas, validações e persistência.
\begin{itemize}
    \item Desenvolvi o CRUD do to-do-list;
    \item Entendi as estrutura de pacotes e importações de classe no Java;
    \item Trabalhei com Repository, Controller, Model, Utils;
    \item Utilizei banco de dados SQL, conceito de ORM.
\end{itemize}
\noindent {\bfseries Tecnologias:} Java, Spring Boot, Maven, JWT, SQL.

%-----------------------------------------    
        
\noindent {\bfseries GPLink} \textbar{}  \href{https://github.com/Lucas-Pontes-Soares/ProjetoTCC}{github.com/Lucas-Pontes-Soares/ProjetoTCC} \hfill\ Fev 2023 -- Out 2023 \\
Plataforma web voltada ao público gamer com o objetivo de centralizar perfis de diferentes redes em um único ecossistema.
\begin{itemize}
    \item Integrei o sistema com APIs e bibliotecas das plataformas Xbox, Steam e PlayStation implementando consumo de dados paginado;
    \item Desenvolvi a interface front-end responsiva para a consolidação e exibição unificada dos perfis e conquistas das redes;
    \item Implementei o fluxo de autenticação e segurança de rotas na aplicação utilizando tokens JWT;
    \item Gerenciei a persistência e a modelagem de dados flexível utilizando o banco de dados NoSQL MongoDB;
    \item Atuei de forma full-stack no design da arquitetura, no desenvolvimento de serviços e na integração com o cliente.
\end{itemize}
\noindent {\bfseries Tecnologias:} React.js, Node.js, MongoDB, API.

%-----------------------------------------    
        
\noindent {\bfseries Upload AI} \textbar{}  \href{https://github.com/Lucas-Pontes-Soares/NLW-IA}{github.com/Lucas-Pontes-Soares/NLW-IA} \hfill\ Set 2023 \\
Plataforma web que consome modelos de IA para gerar automaticamente títulos e descrições com base na transcrição de vídeos.
\begin{itemize}
    \item Desenvolvi a interface web responsiva e modularizada utilizando componentes da biblioteca shadcn/ui;
    \item Implementei a lógica back-end e o pipeline para processamento de áudio e extração da transcrição do vídeo;
    \item Estruturei a engenharia de prompts injetando variáveis de contexto para garantir respostas precisas da IA;
    \item Integrei a API de modelo de linguagem fornecendo a transcrição para a geração automatizada de metadados focados em engajamento.
\end{itemize}
\noindent {\bfseries Tecnologias:} React.js, Node.js, IA, shadcn-ui.

%-----------------------------------------    
        
\noindent {\bfseries La Cafezito} \textbar{}  \href{https://github.com/Lucas-Pontes-Soares/LaCafezito}{github.com/Lucas-Pontes-Soares/LaCafezito} \hfill\ Set 2022 -- Dez 2022 \\
Sistema web para cafeteria com fluxo de montagem de pedidos e gerenciamento interno via painel administrativo com quadro Kanban.
\begin{itemize}
    \item Gerenciei a modelagem do banco de dados relacional e a construção de consultas SQL eficientes utilizando o MySQL;
    \item Desenvolvi as regras de negócio e a lógica de persistência do back-end utilizando a linguagem PHP;
    \item Construí o painel administrativo integrando uma interface Kanban para o controle visual e movimentação dos status dos pedidos.
\end{itemize}
\noindent {\bfseries Tecnologias:} PHP, HTML, CSS, SQL, MySQL.

%-----------------------------------------    
        
\noindent {\bfseries Cancela Automatica} \textbar{}  \href{https://github.com/Lucas-Pontes-Soares/ArduinoCancela}{github.com/Lucas-Pontes-Soares/ArduinoCancela} \hfill\ Fev 2022 -- Out 2022 \\
Sistema de cancela automatizada controlado por microcontrolador, utilizando servomotores, botões e comunicação Bluetooth.
\begin{itemize}
    \item Desenvolvi a lógica computacional em C para o acionamento da cancela simulando a aproximação de composições ferroviárias;
    \item Projetei e montei o circuito físico integrando LEDs de sinalização, sensores e servomotores no microcontrolador;
    \item Implementei o protocolo de comunicação pareado via Bluetooth para viabilizar o controle remoto do sistema por dispositivos móveis.
\end{itemize}
\noindent {\bfseries Tecnologias:} Arduino, C, Sistemas Embarcados.

%-----------------------------------------    
        
\noindent {\bfseries Jogo Tiro ao Alvo} \textbar{}  \href{https://github.com/Lucas-Pontes-Soares/JogoAlvo}{github.com/Lucas-Pontes-Soares/JogoAlvo} \hfill\ Abr 2022 \\
Aplicação web desenvolvida para treinamento de precisão de clique, gerando alvos aleatórios com contagem de acertos e erros.
\begin{itemize}
    \item Desenvolvi a lógica em JavaScript para renderização e posicionamento pseudoaleatório dos alvos na tela;
    \item Implementei um sistema de cronômetro customizável permitindo ao usuário definir o tempo de duração da partida;
    \item Desenvolvi as funções de reset de estado da aplicação para reinicialização instantânea da pontuação e do fluxo do jogo.
\end{itemize}
\noindent {\bfseries Tecnologias:} Front-end, JavaScript, HTML, CSS.

%-----------------------------------------    
        
\noindent {\bfseries Jogo da Forca} \textbar{}  \href{https://github.com/Lucas-Pontes-Soares/JogoForca}{github.com/Lucas-Pontes-Soares/JogoForca} \hfill\ Out 2021 -- Dez 2021 \\
Jogo da forca clássico desenvolvido em C++ com interface via terminal, apresentando escolha de temas e renderização do personagem.
\begin{itemize}
    \item Desenvolvi a lógica do jogo para varredura e identificação de caracteres correspondentes na palavra secreta;
    \item Estruturei a interface visual e a renderização do layout diretamente via console de comandos (terminal);
    \item Implementei o algoritmo de controle de tentativas desenhando o boneco progressivamente a cada entrada incorreta.
\end{itemize}
\noindent {\bfseries Tecnologias:} C++, Lógica.